% Implementation

\section{Implementation}
\label{sec:implementation}

We implement \textsc{WasmSentinel} in approximately 8,500 lines of code across three languages.
This section describes the key implementation details.

\subsection{eBPF Event Collection}

The kernel-side component consists of 450 lines of C code using libbpf.
We attach eBPF programs to the following tracepoints:

\begin{table}[t]
\centering
\caption{eBPF tracepoints for event collection.}
\label{tab:tracepoints}
\begin{tabular}{ll}
\toprule
\textbf{Tracepoint} & \textbf{Event Type} \\
\midrule
\texttt{sys\_enter\_*} & Syscall entry \\
\texttt{sys\_exit\_*} & Syscall return \\
\texttt{sched\_process\_exec} & Process execution \\
\texttt{sched\_process\_exit} & Process termination \\
\texttt{sched\_process\_fork} & Process creation \\
\texttt{tcp\_connect} & Network connections \\
\bottomrule
\end{tabular}
\end{table}

Events are written to a BPF ring buffer (available since Linux 5.8) using the following structure:

\begin{lstlisting}[style=rust, caption={Kernel event structure (C).}, label=lst:event-struct]
struct event {
    __u64 timestamp;
    __u32 pid;
    __u32 ppid;
    __u32 uid;
    __u32 gid;
    __u32 event_type;
    __u32 syscall_nr;
    __s64 ret;
    char comm[16];
    char filename[256];
} __attribute__((packed));
\end{lstlisting}

The ring buffer is sized at 16MB by default, sufficient for approximately 50,000 events before overwriting.
Events are consumed within milliseconds under normal operation, so overflow is rare.

\subsection{Host Program (Go)}

The host application comprises 2,100 lines of Go code and is responsible for:

\textbf{Event Reading.}
We use the \texttt{cilium/ebpf} library to read events from the ring buffer.
A dedicated goroutine polls the buffer every 100$\mu$s, batching events when the arrival rate exceeds 1,000/second.

\textbf{Wasm Engine.}
We embed the Wasmtime runtime~\cite{wasmtime} using cgo bindings.
Wasm modules are compiled ahead-of-time at startup, eliminating JIT overhead during detection.

\textbf{Module Management.}
Each detection rule runs in a separate Wasm instance with:
\begin{itemize}
    \item 64MB maximum linear memory
    \item 100ms execution timeout per event batch
    \item Separate fuel counters for fair scheduling
\end{itemize}

\textbf{Alert Handling.}
Alerts are emitted through a configurable output pipeline supporting:
syslog, Kafka, file output, and HTTP webhooks.

\subsection{Detection Module (Rust → Wasm)}

The ML detection module comprises 1,850 lines of Rust code, organized as follows:

\begin{lstlisting}[style=rust, caption={Module structure.}]
ml-detector/
  src/
    lib.rs       // Wasm exports, entry point
    stide.rs     // STIDE n-gram detector
    mlp.rs       // MLP neural network
    hybrid.rs    // Hybrid scheduler
    features.rs  // Feature extraction
\end{lstlisting}

\textbf{Wasm Exports.}
The module exports four primary functions:

\begin{lstlisting}[style=rust, caption={Wasm exported functions.}]
#[no_mangle]
pub extern "C" fn init(cfg: *const u8, len: usize) -> i32;

#[no_mangle]
pub extern "C" fn detect(evt: *const u8, len: usize) -> i32;

#[no_mangle]
pub extern "C" fn train(data: *const u8, len: usize) -> i32;

#[no_mangle]
pub extern "C" fn get_stats(len: *mut usize) -> *const u8;
\end{lstlisting}

\textbf{Dependencies.}
We minimize dependencies to reduce binary size:
\begin{itemize}
    \item \texttt{serde\_json}: Event parsing (adds 40KB)
    \item No external ML libraries
    \item No standard library allocator (using \texttt{wee\_alloc})
\end{itemize}

The compiled Wasm module is 180KB, compared to 10--50MB for typical ML deployments.

\subsection{Model Training (Python)}

We train the MLP model using a pure-NumPy implementation (no TensorFlow/PyTorch dependency) in 650 lines of Python:

\textbf{Data Loading.}
Training data is loaded from the ADFA-LD dataset, which contains 5,951 syscall traces:
\begin{itemize}
    \item Normal traces: 5,205 (87.5\%)
    \item Attack traces: 746 (12.5\%)
\end{itemize}

\textbf{Class Balancing.}
To address class imbalance, we apply weighted cross-entropy loss:
\begin{equation}
\mathcal{L} = -\frac{1}{N}\sum_{i=1}^{N} w_{y_i} \cdot [y_i \log(\hat{y}_i) + (1-y_i)\log(1-\hat{y}_i)]
\end{equation}
where $w_1 = 3.87$ (attack) and $w_0 = 0.57$ (normal) are inversely proportional to class frequencies.

\textbf{Weight Export.}
After training, weights are exported to JSON format for loading into the Wasm module:

\begin{lstlisting}[style=rust, caption={Weight file format (JSON).}]
[
  {
    "weights": [[...], ...],  // 64x128 matrix
    "biases": [...]           // 128 vector
  },
  // ... additional layers
]
\end{lstlisting}

The exported weight file is 558KB for the 4-layer architecture.

\subsection{Build and Deployment}

\textbf{Compilation.}
The Wasm module is compiled using:
\begin{verbatim}
cargo build --target wasm32-wasip1 --release
\end{verbatim}

With release optimizations (\texttt{lto = true, opt-level = "s"}), the module compiles in approximately 30 seconds.

\textbf{Deployment.}
\textsc{WasmSentinel} is distributed as a single binary with embedded eBPF bytecode and Wasm modules.
Deployment requires:
\begin{itemize}
    \item Linux kernel 5.8+ (for ring buffer support)
    \item Root privileges (for eBPF loading)
    \item No additional runtime dependencies
\end{itemize}
